% 第六章 实验评估
\section{实验评估}

\subsection{实验环境}

\begin{table}[H]
    \centering
    \caption{实验环境配置}
    \label{tab:env}
    \begin{tabularx}{\textwidth}{lX}
        \toprule
        \textbf{项目} & \textbf{配置} \\
        \midrule
        处理器 & 12th Gen Intel Core i7-12800HX（24核） \\
        内存 & 15.75 GB \\
        操作系统 & Windows\_NT 10.0.26200 \\
        Node.js 版本 & v20.19.4（V8 11.3.244.8） \\
        椭圆曲线库 & elliptic.js v6.5.4（secp256k1） \\
        \bottomrule
    \end{tabularx}
\end{table}

\subsection{功能正确性验证}

\subsubsection{Shamir 秘密共享正确性}

为验证 Shamir 秘密共享实现的正确性，我们在 $n \in \{3, 5, 7, 10\}$、$t \in \{2, 3, 5, 7\}$ 的参数组合下各进行 100 轮随机测试。测试过程为：随机生成 256 位主密钥 $\rightarrow$ 使用 $t$-of-$n$ 拆分 $\rightarrow$ 随机选取 $t$ 个份额 $\rightarrow$ Lagrange 插值恢复 $\rightarrow$ 比较与原始主密钥的一致性。

\begin{table}[H]
    \centering
    \caption{Shamir 正确性验证结果}
    \label{tab:shamir_test}
    \begin{tabularx}{\textwidth}{cccccc}
        \toprule
        $n$ & $t$ & 测试轮数 & 成功次数 & 正确率 & 平均耗时(ms) \\
        \midrule
        3 & 2 & 100 & 100 & 100.00\% & 0.0101 \\
        5 & 3 & 100 & 100 & 100.00\% & 0.0123 \\
        7 & 5 & 100 & 100 & 100.00\% & 0.0629 \\
        10 & 7 & 100 & 100 & 100.00\% & 0.0406 \\
        \bottomrule
    \end{tabularx}
\end{table}

\subsubsection{同态刷新正确性验证}

验证刷新后的份额使用 Lagrange 插值恢复的主密钥与原始主密钥一致。测试参数同上。

\textbf{验证步骤}：（1）生成主密钥 $s$；（2）拆分得到 $v_i, r_i$；（3）生成零和多项式刷新得到 $\delta_{v,i}, \delta_{r,i}$；（4）计算新份额 $v'_i = v_i + \delta_{v,i}$；（5）用 $t$ 个 $v'_i$ 恢复 $s'$；（6）验证 $s' = s$。

\begin{table}[H]
    \centering
    \caption{同态刷新正确性验证结果}
    \label{tab:refresh_test}
    \begin{tabularx}{\textwidth}{cccccc}
        \toprule
        $n$ & $t$ & 测试轮数 & 刷新前正确率 & 刷新后正确率 & 刷新耗时(ms) \\
        \midrule
        3 & 2 & 100 & 100.00\% & 100.00\% & 0.0304 \\
        5 & 3 & 100 & 100.00\% & 100.00\% & 0.0671 \\
        7 & 5 & 100 & 100.00\% & 100.00\% & 0.1746 \\
        10 & 7 & 100 & 100.00\% & 100.00\% & 0.2480 \\
        \bottomrule
    \end{tabularx}
\end{table}

\subsection{密码学运算性能测试}

\begin{table}[H]
    \centering
    \caption{各密码学操作的性能开销}
    \label{tab:perf}
    \begin{tabularx}{\textwidth}{lc}
        \toprule
        \textbf{操作} & \textbf{平均耗时(ms)} \\
        \midrule
        主密钥生成（256 位随机数） & 0.0206 \\
        Pedersen 双多项式拆分（5-of-3） & 0.0940 \\
        Pedersen 承诺生成 & 1.3569 \\
        零和多项式增量生成（5 份） & 0.0693 \\
        承诺同态刷新计算（5 份） & 6.8089 \\
        Lagrange 插值恢复主密钥（3 份） & 0.0054 \\
        ECIES 加密（份额值） & 1.5716 \\
        主承诺同态验证（3 份） & 1.6624 \\
        \bottomrule
    \end{tabularx}
\end{table}

\subsection{规模扩展性测试}

测试不同监护人数目 $n$ 和门限 $t$ 下的性能变化趋势（50轮测试取平均）。

\begin{table}[H]
    \centering
    \caption{规模扩展性测试结果}
    \label{tab:scalability}
    \begin{tabularx}{\textwidth}{cccc}
        \toprule
        $(n, t)$ & 拆分耗时(ms) & 刷新耗时(ms) & 恢复耗时(ms) \\
        \midrule
        (3, 2) & 0.054 & 0.034 & 0.004 \\
        (5, 3) & 0.098 & 0.125 & 0.007 \\
        (7, 5) & 0.175 & 0.154 & 0.016 \\
        (10, 7) & 0.299 & 0.239 & 0.032 \\
        \bottomrule
    \end{tabularx}
\end{table}

\subsection{同态验证精度测试}

测试主承诺同态验证的精度（准确率、召回率）：

\begin{table}[H]
    \centering
    \caption{同态验证精度测试}
    \label{tab:verify}
    \begin{tabularx}{\textwidth}{lccccc}
        \toprule
        \textbf{指标} & \textbf{TP} & \textbf{FN} & \textbf{FP} & \textbf{TN} & \textbf{Accuracy} \\
        \midrule
        综合（100轮） & 100 & 0 & 0 & 100 & 100.00\% \\
        \bottomrule
    \end{tabularx}
\end{table}

\subsection{与现有方案的对比测试}

\begin{table}[H]
    \centering
    \caption{与现有方案的对比}
    \label{tab:compare_tools}
    \begin{tabularx}{\textwidth}{lccc}
        \toprule
        \textbf{特性} & \textbf{本方案} & \textbf{Ledger Recover} & \textbf{Safe\{Wallet\}} \\
        \midrule
        零信任 Server & $\checkmark$ & $\times$ & $\checkmark$ \\
        非链上资产支持 & $\checkmark$ & $\times$ & $\times$ \\
        份额刷新 & $\checkmark$（同态） & $\times$ & $\times$ \\
        抗胁迫 & $\checkmark$ & $\times$ & $\times$ \\
        隐私保护 & 全流程加密 & KYC 可见 & 链上公开 \\
        开源 & 全栈开源 & 闭源 & 部分 \\
        \bottomrule
    \end{tabularx}
\end{table}

\textbf{实验结论}：本文提出的数字遗产管家系统在功能正确性上通过了全部 100 轮测试，Shamir 秘密共享拆分与恢复、同态份额刷新、主承诺同态验证的正确率均达到 100\%。密码学操作性能方面，核心操作（承诺生成、同态验证、ECIES 加解密）的平均耗时均在 0.9--6.8ms 范围内，满足实时交互的需求。规模扩展性测试表明，各操作的耗时随监护人数目呈线性增长趋势，10-of-7 配置下单次拆分仅需 0.3ms，系统具备良好的可扩展性。与现有方案（Ledger Recover、Safe\{Wallet\}）的对比显示，本方案在零信任 Server、非链上资产支持、份额刷新和抗胁迫设计方面具有显著优势。
